\section{Les schémas de Chevalley}
\label{section:I.8-fr}

\subsection{Anneaux locaux apparentés.}
\label{subsection:I.8.1-fr}

Pour tout anneau local $A$, nous noterons $\mathfrak m(A)$ l'idéal maximal de
$A$.

\begin{lemma}[8.1.1]
\label{I.8.1.1-fr}
Soient $A$, $B$ deux anneaux locaux tels que $A\subset B$~; les conditions
suivantes sont équivalentes~:
\begin{enumerate}
  \item[(i)] $\mathfrak m(B)\cap A=\mathfrak m(A)$~;
  \item[(ii)] $\mathfrak m(A)\subset\mathfrak m(B)$~;
  \item[(iii)] $1$ n'appartient pas à l'idéal de $B$ engendré par
    $\mathfrak m(A)$.
\end{enumerate}
\end{lemma}

Il est évident que (i) entraîne (ii) et (ii) entraîne (iii)~; enfin, si (iii)
est vérifiée, $\mathfrak m(B)\cap A$ contient $\mathfrak m(A)$ et ne contient
pas $1$, donc est égal à $\mathfrak m(A)$.

Lorsque les conditions équivalentes de
(\hyperref[I.8.1.1-fr]{8.1.1}) sont remplies, on dit que $B$ \emph{domine}
$A$~; il revient au même de dire que l'injection $A\to B$ est un
homomorphisme \emph{local}. Il est clair que, dans l'ensemble des sous-anneaux
locaux d'un anneau $R$, la relation de domination est une relation d'ordre.

\begin{env}[8.1.2]
\label{I.8.1.2-fr}
Considérons maintenant un \emph{corps} $R$. Pour tout sous-anneau $A$ de $R$,
nous désignerons par $L(A)$ l'ensemble des anneaux locaux $A_\mathfrak p$, où
$\mathfrak p$ parcourt le spectre premier de $A$~; ils sont identifiés à des
sous-anneaux de $R$ contenant $A$. Comme
$\mathfrak p=(\mathfrak p A_\mathfrak p)\cap A$, l'application
$\mathfrak p\mapsto A_\mathfrak p$ de $\operatorname{Spec}(A)$ dans $L(A)$
est bijective.
\end{env}

\begin{lemma}[8.1.3]
\label{I.8.1.3-fr}
Soient $R$ un corps, $A$ un sous-anneau de $R$. Pour qu'un sous-anneau local
$M$ de $R$ domine un anneau $A_\mathfrak p\in L(A)$, il faut et il suffit que
$A\subset M$~; l'anneau local $A_\mathfrak p$ dominé par $M$ est alors unique
et correspond à $\mathfrak p=\mathfrak m(M)\cap A$.
\end{lemma}

En effet, si $M$ domine $A_\mathfrak p$, on a
$\mathfrak m(M)\cap A_\mathfrak p=\mathfrak p A_\mathfrak p$ d'après
(\hyperref[I.8.1.1-fr]{8.1.1}), d'où l'unicité de $\mathfrak p$~; d'autre
part, si $A\subset M$, $\mathfrak m(M)\cap A=\mathfrak p$ est premier dans
$A$, et comme $A-\mathfrak p\subset M$, on a $A_\mathfrak p\subset M$ et
$\mathfrak p A_\mathfrak p\subset\mathfrak m(M)$, donc $M$ domine
$A_\mathfrak p$.

\oldpage[I]{165}
\begin{lemma}[8.1.4]
\label{I.8.1.4-fr}
Soient $R$ un corps, $M$, $N$ deux sous-anneaux locaux de $R$, $P$ le
sous-anneau de $R$ engendré par $M\cup N$. Les conditions suivantes sont
équivalentes~:
\begin{enumerate}
  \item[(i)] Il existe un idéal premier $\mathfrak p$ de $P$ tel que
    $\mathfrak m(M)=\mathfrak p\cap M$,
    $\mathfrak m(N)=\mathfrak p\cap N$.
  \item[(ii)] L'idéal $\mathfrak a$ engendré dans $P$ par
    $\mathfrak m(M)\cup\mathfrak m(N)$ est distinct de $P$.
  \item[(iii)] Il existe un sous-anneau local $Q$ de $R$ dominant à la fois
    $M$ et $N$.
\end{enumerate}
\end{lemma}

Il est clair que (i) implique (ii)~; inversement, si $\mathfrak a\ne P$,
$\mathfrak a$ est contenu dans un idéal maximal $\mathfrak n$ de $P$ et comme
$1\notin\mathfrak n$, $\mathfrak n\cap M$ contient $\mathfrak m(M)$ et est
distinct de $M$, donc $\mathfrak n\cap M=\mathfrak m(M)$ et de même
$\mathfrak n\cap N=\mathfrak m(N)$. Il est clair que si $Q$ domine $M$ et
$N$, on a $P\subset Q$, et
\[
  \mathfrak m(M)=\mathfrak m(Q)\cap M
  =(\mathfrak m(Q)\cap P)\cap M,\qquad
  \mathfrak m(N)=(\mathfrak m(Q)\cap P)\cap N,
\]
donc (iii) entraîne (i)~; la réciproque est évidente en prenant
$Q=P_\mathfrak p$.

Lorsque les conditions de (\hyperref[I.8.1.4-fr]{8.1.4}) sont satisfaites,
on dit, avec C.~Chevalley, que les anneaux locaux $M$ et $N$ sont
\emph{apparentés}.

\begin{proposition}[8.1.5]
\label{I.8.1.5-fr}
Soient $A$, $B$ deux sous-anneaux d'un corps $R$, $C$ le sous-anneau de $R$
engendré par $A\cup B$. Les conditions suivantes sont équivalentes~:
\begin{enumerate}
  \item[(i)] Pour tout anneau local $Q$ contenant $A$ et $B$, on a
    $A_\mathfrak p=B_\mathfrak q$, en posant
    $\mathfrak p=\mathfrak m(Q)\cap A$,
    $\mathfrak q=\mathfrak m(Q)\cap B$.
  \item[(ii)] Pour tout idéal premier $\mathfrak r$ de $C$, on a
    $A_\mathfrak p=B_\mathfrak q$, en posant
    $\mathfrak p=\mathfrak r\cap A$, $\mathfrak q=\mathfrak r\cap B$.
  \item[(iii)] Si $M\in L(A)$ et $N\in L(B)$ sont apparentés, ils sont
    identiques.
  \item[(iv)] On a $L(A)\cap L(B)=L(C)$.
\end{enumerate}
\end{proposition}

Les lemmes (\hyperref[I.8.1.3-fr]{8.1.3}) et
(\hyperref[I.8.1.4-fr]{8.1.4}) prouvent que (i) et (iii) sont équivalentes~;
il est clair que (i) entraîne (ii) en l'appliquant à $Q=C_\mathfrak r$~;
inversement, (ii) entraîne (i), car si $Q$ contient $A\cup B$, il contient
$C$, et si $\mathfrak r=\mathfrak m(Q)\cap C$, on a
$\mathfrak p=\mathfrak r\cap A$ et $\mathfrak q=\mathfrak r\cap B$ d'après
(\hyperref[I.8.1.3-fr]{8.1.3}). Il est immédiat que (iv) implique (i), car si
$Q$ contient $A\cup B$, il domine un anneau local
$C_\mathfrak r\in L(C)$ par (\hyperref[I.8.1.3-fr]{8.1.3})~; on a par
hypothèse $C_\mathfrak r\in L(A)\cap L(B)$, et
(\hyperref[I.8.1.1-fr]{8.1.1}) et (\hyperref[I.8.1.3-fr]{8.1.3}) prouvent que
$C_\mathfrak r=A_\mathfrak p=B_\mathfrak q$. Prouvons enfin que (iii)
entraîne (iv). Soit $Q\in L(C)$~; $Q$ domine un $M\in L(A)$ et un
$N\in L(B)$ (\hyperref[I.8.1.3-fr]{8.1.3}), donc $M$ et $N$, étant
apparentés, sont identiques par hypothèse. Comme on a alors $C\subset M$,
$M$ domine un $Q'\in L(C)$ (\hyperref[I.8.1.3-fr]{8.1.3}), donc $Q$ domine
$Q'$, ce qui (\hyperref[I.8.1.3-fr]{8.1.3}) entraîne nécessairement
$Q=Q'=M$, donc $Q\in L(A)\cap L(B)$. Inversement, si
$Q\in L(A)\cap L(B)$, on a $C\subset Q$, donc
(\hyperref[I.8.1.3-fr]{8.1.3}) $Q$ domine un
$Q''\in L(C)\subset L(A)\cap L(B)$~; $Q$ et $Q''$ étant apparentés sont
identiques, donc $Q''=Q\in L(C)$, ce qui achève la démonstration.

\subsection{Anneaux locaux d'un schéma intègre.}
\label{subsection:I.8.2-fr}

\begin{env}[8.2.1]
\label{I.8.2.1-fr}
Soient $X$ un préschéma \emph{intègre}, $R$ son corps des fonctions
rationnelles, identique à l'anneau local du point générique $a$ de $X$~; pour
tout $x\in X$, on sait que $\mathscr{O}_x$ s'identifie canoniquement à un
sous-anneau de $R$ (\hyperref[I.7.1.5-fr]{7.1.5}), et pour toute fonction
rationnelle $f\in R$, le domaine de définition $\delta(f)$ de $f$ est
l'ensemble ouvert des $x\in X$ tels que $f\in\mathscr{O}_x$. Il en résulte
(\hyperref[I.7.2.6-fr]{7.2.6}) que pour tout ouvert $U\subset X$, on a
\begin{equation}
  \Gamma(U,\mathscr{O}_X)=\bigcap_{x\in U}\mathscr{O}_x.
  \tag{8.2.1.1}\label{I.8.2.1.1-fr}
\end{equation}
\end{env}

\oldpage[I]{166}
\begin{proposition}[8.2.2]
\label{I.8.2.2-fr}
Soient $X$ un préschéma intègre, $R$ son corps des fonctions rationnelles.
Pour que $X$ soit un schéma, il faut et il suffit que la relation
«~$\mathscr{O}_x$ et $\mathscr{O}_y$ sont apparentés~»
(\hyperref[I.8.1.4-fr]{8.1.4}) entre points $x$, $y$ de $X$ implique $x=y$.
\end{proposition}

Supposons cette condition vérifiée, et montrons que $X$ est séparé. Soient
$U$ et $V$ deux ouverts affines distincts de $X$, $A$ et $B$ leurs anneaux,
identifiés à des sous-anneaux de $R$~; $U$ (resp. $V$) s'identifie donc
(\hyperref[I.8.1.2-fr]{8.1.2}) à $L(A)$ (resp. $L(B)$), et l'hypothèse
entraîne (\hyperref[I.8.1.5-fr]{8.1.5}) que si $C$ est le sous-anneau de $R$
engendré par $A\cup B$, $W=U\cap V$ s'identifie à
$L(A)\cap L(B)=L(C)$. En outre, on sait
(\cite{I-1}, p.~5-03, prop.~4 \emph{bis}) que tout sous-anneau $E$ de $R$ est
égal à l'intersection des anneaux locaux appartenant à $L(E)$~; $C$
s'identifie donc à l'intersection des anneaux $\mathscr{O}_z$ pour $z\in W$,
autrement dit (\hyperref[I.8.2.1.1-fr]{8.2.1.1}) à
$\Gamma(W,\mathscr{O}_X)$. Considérons alors le sous-préschéma induit par $X$
sur $W$~; à l'homomorphisme identique
$\varphi:C\to\Gamma(W,\mathscr{O}_X)$ correspond
(\hyperref[I.2.2.4-fr]{2.2.4}) un morphisme
$\Phi=(\psi,\theta):W\to\operatorname{Spec}(C)$~; nous allons voir que $\Phi$
est un \emph{isomorphisme} de préschémas, d'où résultera que $W$ est un ouvert
\emph{affine}. L'identification de $W$ à
$L(C)=\operatorname{Spec}(C)$ montre que $\psi$ est \emph{bijective}. D'autre
part, pour tout $x\in W$, $\theta_x^\sharp$ est l'injection
$C_\mathfrak r\to\mathscr{O}_x$, si
$\mathfrak r=\mathfrak m_x\cap C$, et par définition $C_\mathfrak r$ est
identifié à $\mathscr{O}_x$, donc $\theta_x^\sharp$ est bijective. Il reste
donc à voir que $\psi$ est un \emph{homéomorphisme}, autrement dit que pour
toute partie fermée $F\subset W$, $\psi(F)$ est fermé dans
$\operatorname{Spec}(C)$. Or, $F$ est la trace sur $W$ d'une partie fermée de
$U$, de la forme $V(\mathfrak a)$, où $\mathfrak a$ est un idéal de $A$~;
montrons que $\psi(F)=V(\mathfrak a C)$, ce qui prouvera notre assertion. En
effet, les idéaux premiers de $C$ contenant $\mathfrak a C$ sont les idéaux
premiers de $C$ contenant $\mathfrak a$, donc les idéaux de la forme
$\psi(x)=\mathfrak m_x\cap C$ où
$\mathfrak a\subset\mathfrak m_x$ et $x\in W$~; comme
$\mathfrak a\subset\mathfrak m_x$ équivaut à
$x\in V(\mathfrak a)=W\cap F$ pour $x\in U$, on a bien
$\psi(F)=V(\mathfrak a C)$.

Il s'ensuit que $X$ est séparé, car $U\cap V$ est affine et son anneau $C$ est
engendré par la réunion $A\cup B$ des anneaux de $U$ et $V$
(\hyperref[I.5.5.6-fr]{5.5.6}).

Inversement, supposons $X$ séparé, et soient $x$, $y$ deux points de $X$ tels
que $\mathscr{O}_x$ et $\mathscr{O}_y$ soient apparentés. Soit $U$ (resp.
$V$) un ouvert affine contenant $x$ (resp. $y$), d'anneau $A$ (resp. $B$)~;
on sait alors que $U\cap V$ est affine et que son anneau $C$ est engendré par
$A\cup B$ (\hyperref[I.5.5.6-fr]{5.5.6}). Si
$\mathfrak p=\mathfrak m_x\cap A$, $\mathfrak q=\mathfrak m_y\cap B$, on a
$A_\mathfrak p=\mathscr{O}_x$, $B_\mathfrak q=\mathscr{O}_y$, et comme
$A_\mathfrak p$ et $B_\mathfrak q$ sont apparentés, il existe un idéal premier
$\mathfrak r$ de $C$ tel que $\mathfrak p=\mathfrak r\cap A$,
$\mathfrak q=\mathfrak r\cap B$ (\hyperref[I.8.1.4-fr]{8.1.4}). Mais alors il
existe un point $z\in U\cap V$ tel que
$\mathfrak r=\mathfrak m_z\cap C$ puisque $U\cap V$ est affine, et on a
évidemment $x=z$, et $y=z$, d'où $x=y$.

\begin{corollary}[8.2.3]
\label{I.8.2.3-fr}
Soient $X$ un schéma intègre, $x$, $y$ deux points de $X$. Pour que
$x\in\overline{\{y\}}$, il faut et il suffit que
$\mathscr{O}_x\subset\mathscr{O}_y$, autrement dit que toute fonction
rationnelle définie en $x$ soit définie en $y$.
\end{corollary}

La condition est évidemment nécessaire puisque le domaine de définition
$\delta(f)$ d'une fonction rationnelle $f\in R$ est ouvert~; montrons qu'elle
est suffisante. Si $\mathscr{O}_x\subset\mathscr{O}_y$, il existe un idéal
premier $\mathfrak p$ de $\mathscr{O}_x$ tel que $\mathscr{O}_y$ domine
$(\mathscr{O}_x)_\mathfrak p$ (\hyperref[I.8.1.3-fr]{8.1.3})~; or
(\hyperref[I.2.4.2-fr]{2.4.2}) il existe $z\in X$ tel que
$x\in\overline{\{z\}}$ et que
$\mathscr{O}_z=(\mathscr{O}_x)_\mathfrak p$~; comme $\mathscr{O}_z$ et
$\mathscr{O}_y$ sont apparentés, on a $z=y$ par
(\hyperref[I.8.2.2-fr]{8.2.2}), d'où le corollaire.

\begin{corollary}[8.2.4]
\label{I.8.2.4-fr}
Si $X$ est un schéma intègre, l'application $x\mapsto\mathscr{O}_x$ est
injective~; autrement dit, si $x$, $y$ sont deux points distincts de $X$, il
existe une fonction rationnelle définie en l'un de ces points et non en
l'autre.
\end{corollary}

\oldpage[I]{167}
Cela résulte de (\hyperref[I.8.2.3-fr]{8.2.3}) et de l'axiome
$(T_0)$ (\hyperref[I.2.1.4-fr]{2.1.4}).

\begin{corollary}[8.2.5]
\label{I.8.2.5-fr}
Soit $X$ un schéma intègre dont l'espace sous-jacent est noethérien~;
lorsque $f$ parcourt le corps $R$ des fonctions rationnelles sur $X$, les
ensembles $\delta(f)$ engendrent la topologie de $X$.
\end{corollary}

En effet, toute partie fermée de $X$ est alors réunion finie d'ensembles
fermés irréductibles, c'est-à-dire de la forme $\overline{\{y\}}$
(\hyperref[I.2.1.5-fr]{2.1.5}). Or, si
$x\notin\overline{\{y\}}$, il existe une fonction rationnelle $f$ définie en
$x$ et non en $y$ (\hyperref[I.8.2.3-fr]{8.2.3}), autrement dit, on a
$x\in\delta(f)$ et $\delta(f)$ ne rencontre pas $\overline{\{y\}}$. Le
complémentaire de $\overline{\{y\}}$ est par suite réunion d'ensembles de la
forme $\delta(f)$, et en vertu de la première remarque, tout ouvert de $X$ est
réunion d'intersections finies d'ouverts de la forme $\delta(f)$.

\begin{env}[8.2.6]
\label{I.8.2.6-fr}
Le cor. (\hyperref[I.8.2.5-fr]{8.2.5}) montre que la topologie de $X$ est
entièrement caractérisée par la donnée de la famille d'anneaux locaux
$(\mathscr{O}_x)_{x\in X}$ ayant $R$ pour corps des fractions. Il revient au
même d'ailleurs de dire que les parties fermées de $X$ sont définies de la
façon suivante~: étant donnée une partie finie
$\{x_1,\ldots,x_n\}$ de $X$, on considère l'ensemble des $y\in X$ tels que
$\mathscr{O}_y\subset\mathscr{O}_{x_i}$ pour un indice $i$ au moins, et ces
ensembles (pour tous les choix de $\{x_1,\ldots,x_n\}$) sont les ensembles
fermés de $X$. En outre, une fois connue la topologie de $X$, le faisceau
structural $\mathscr{O}_X$ est aussi bien déterminé par la famille des
$\mathscr{O}_x$, puisque
$\Gamma(U,\mathscr{O}_X)=\bigcap_{x\in U}\mathscr{O}_x$
(\hyperref[I.8.2.1.1-fr]{8.2.1.1}). La famille
$(\mathscr{O}_x)_{x\in X}$ détermine donc complètement le préschéma $X$
lorsque $X$ est un schéma intègre dont l'espace sous-jacent est noethérien.
\end{env}

\begin{proposition}[8.2.7]
\label{I.8.2.7-fr}
Soient $X$, $Y$ deux schémas intègres, $f:X\to Y$ un morphisme dominant
(\hyperref[I.2.2.6-fr]{2.2.6}), $K$ (resp. $L$) le corps des fonctions
rationnelles sur $X$ (resp. $Y$). Alors $L$ s'identifie à un sous-corps de
$K$, et pour tout $x\in X$, $\mathscr{O}_{f(x)}$ est l'unique anneau local de
$Y$ dominé par $\mathscr{O}_x$.
\end{proposition}

En effet, si $f=(\psi,\theta)$ et si $a$ est le point générique de $X$,
$\psi(a)$ est le point générique de $Y$
(\hyperref[0.2.1.5-fr]{0, 2.1.5})~; $\theta_a^\sharp$ est par suite un
monomorphisme du corps $L=\mathscr{O}_{\psi(a)}$ dans le corps
$K=\mathscr{O}_a$. Comme tout ouvert affine non vide $U$ de $Y$ contient
$\psi(a)$, il résulte de (\hyperref[I.2.2.4-fr]{2.2.4}) que l'homomorphisme
$\Gamma(U,\mathscr{O}_Y)\to\Gamma(\psi^{-1}(U),\mathscr{O}_X)$ correspondant
à $f$ est la restriction de $\theta_a^\sharp$ à
$\Gamma(U,\mathscr{O}_Y)$. Donc, pour tout $x\in X$,
$\theta_x^\sharp$ est la restriction à $\mathscr{O}_{\psi(x)}$ de
$\theta_a^\sharp$, et est par suite un monomorphisme. On sait en outre que
$\theta_x^\sharp$ est un homomorphisme local, donc, si on identifie $L$ à un
sous-corps de $K$ par $\theta_a^\sharp$, $\mathscr{O}_{\psi(x)}$ est dominé
par $\mathscr{O}_x$ (\hyperref[I.8.1.1-fr]{8.1.1})~; c'est d'ailleurs le seul
anneau local de $Y$ dominé par $\mathscr{O}_x$ puisque deux anneaux locaux de
$Y$ qui sont apparentés sont identiques
(\hyperref[I.8.2.2-fr]{8.2.2}).

\begin{proposition}[8.2.8]
\label{I.8.2.8-fr}
Soient $X$ un préschéma irréductible, $f:X\to Y$ une immersion locale
(resp. un isomorphisme local)~; on suppose en outre le morphisme $f$ séparé.
Alors $f$ est une immersion (resp. une immersion ouverte).
\end{proposition}

Soit $f=(\psi,\theta)$~; il suffit, dans les deux cas, de prouver que $\psi$
est un \emph{homéomorphisme} de $X$ sur $\psi(X)$
(\hyperref[I.4.5.3-fr]{4.5.3}). Remplaçant $f$ par $f_{\mathrm{red}}$
(\hyperref[I.5.1.6-fr]{5.1.6} et
\hyperref[I.5.5.1-fr]{5.5.1, (vi)}), on peut supposer $X$ et $Y$ réduits. Si
$Y'$ est le sous-préschéma fermé réduit de $Y$ ayant pour espace sous-jacent
$\overline{\psi(X)}$, $f$ se factorise en
$X\xrightarrow{f'}Y'\xrightarrow{j}Y$, où $j$ est l'injection canonique
(\hyperref[I.5.2.2-fr]{5.2.2}). Il résulte de
(\hyperref[I.5.5.1-fr]{5.5.1, (v)}) que $f'$ est encore un morphisme séparé~;
en outre, $f'$ est encore

\oldpage[I]{168}
une immersion locale (resp. un isomorphisme local), car la question étant
locale sur $X$ et $Y$, on peut se borner pour le démontrer au cas où $f$ est
une immersion fermée (resp. une immersion ouverte) et notre assertion découle
alors aussitôt de (\hyperref[I.4.2.2-fr]{4.2.2}).

On peut donc supposer que $f$ est un morphisme \emph{dominant}, ce qui
entraîne que $Y$ est, lui aussi, irréductible
(\hyperref[0.2.1.5-fr]{0, 2.1.5}), donc que $X$ et $Y$ sont tous deux
\emph{intègres}. En outre, la question étant locale sur $Y$, on peut supposer
que $Y$ est un schéma affine~; comme $f$ est séparé, $X$ est un schéma
(\hyperref[I.5.5.1-fr]{5.5.1, (ii)}), et on est finalement dans les
hypothèses de (\hyperref[I.8.2.7-fr]{8.2.7}). Alors, pour tout $x\in X$,
$\theta_x^\sharp$ est injectif~; mais l'hypothèse que $f$ est une immersion
locale implique que $\theta_x^\sharp$ est surjectif
(\hyperref[I.4.2.2-fr]{4.2.2}), donc $\theta_x^\sharp$ est bijectif,
autrement dit (avec l'identification de
(\hyperref[I.8.2.7-fr]{8.2.7})) on a
$\mathscr{O}_{\psi(x)}=\mathscr{O}_x$. Cela implique par
(\hyperref[I.8.2.4-fr]{8.2.4}) que $\psi$ est une application
\emph{injective}, ce qui démontre déjà la proposition lorsque $f$ est un
isomorphisme local (\hyperref[I.4.5.3-fr]{4.5.3}). Lorsqu'on suppose seulement
que $f$ est une immersion locale, pour tout $x\in X$ il existe un voisinage
ouvert $U$ de $x$ dans $X$ et un voisinage ouvert $V$ de $\psi(x)$ dans $Y$
tels que la restriction de $\psi$ à $U$ soit un homéomorphisme de $U$ sur une
partie \emph{fermée} de $V$. Or, $U$ est dense dans $X$, donc $\psi(U)$ est
dense dans $Y$ et \emph{a fortiori} dans $V$, ce qui prouve que
$\psi(U)=V$~; comme $\psi$ est injective, $\psi^{-1}(V)=U$ et ceci achève de
montrer que $\psi$ est un homéomorphisme de $X$ sur $\psi(X)$.

\subsection{Les schémas de Chevalley.}
\label{subsection:I.8.3-fr}

\begin{env}[8.3.1]
\label{I.8.3.1-fr}
Soient $X$ un schéma intègre \emph{noethérien}, $R$ son corps des fonctions
rationnelles~; désignons par $X'$ l'ensemble des sous-anneaux locaux
$\mathscr{O}_x\subset R$, où $x$ parcourt $X$. L'ensemble $X'$ vérifie les
trois conditions suivantes~:
\begin{enumerate}
  \item[(Sch. 1)] \emph{Pour tout $M\in X'$, $R$ est le corps des fractions
    de $M$.}
  \item[(Sch. 2)] \emph{Il existe un ensemble fini de sous-anneaux
    noethériens $A_i$ de $R$ tels que $X'=\bigcup_i L(A_i)$ et que, pour tout
    couple d'indices $i$, $j$, le sous-anneau $A_{ij}$ de $R$ engendré par
    $A_i\cup A_j$ soit une algèbre de type fini sur $A_i$.}
  \item[(Sch. 3)] \emph{Deux éléments $M$, $N$ de $X'$ qui sont apparentés
    sont identiques.}
\end{enumerate}
\end{env}

On a en effet vu dans (\hyperref[I.8.2.1-fr]{8.2.1}) que (Sch. 1) est
satisfaite, et (Sch. 3) résulte de (\hyperref[I.8.2.2-fr]{8.2.2}). Pour
démontrer (Sch. 2), il suffit de recouvrir $X$ par un nombre fini d'ouverts
affines $U_i$, d'anneaux noethériens, et de prendre
$A_i=\Gamma(U_i,\mathscr{O}_X)$~; l'hypothèse que $X$ est un schéma entraîne
que $U_i\cap U_j$ est affine et que
$\Gamma(U_i\cap U_j,\mathscr{O}_X)=A_{ij}$
(\hyperref[I.5.5.6-fr]{5.5.6})~; en outre, comme l'espace $U_i$ est
noethérien, l'immersion $U_i\cap U_j\to U_i$ est de type fini
(\hyperref[I.6.3.5-fr]{6.3.5}), donc $A_{ij}$ est une $A_i$-algèbre de type
fini (\hyperref[I.6.3.3-fr]{6.3.3}).

\begin{env}[8.3.2]
\label{I.8.3.2-fr}
Les structures dont les axiomes sont (Sch. 1), (Sch. 2) et (Sch. 3)
généralisent les «~schémas~» au sens de C.~Chevalley, qui suppose en outre que
$R$ est une extension de type fini d'un corps $K$ et que les $A_i$ sont des
$K$-algèbres de type fini (ce qui rend inutile une partie de (Sch. 2))
\cite{I-1}. Inversement, si on a une telle structure sur un ensemble $X'$,
on peut lui associer un schéma intègre $X$ en utilisant les remarques de
(\hyperref[I.8.2.6-fr]{8.2.6})~: l'espace sous-jacent de $X$ est égal à $X'$
muni de la topologie définie dans (\hyperref[I.8.2.6-fr]{8.2.6}), et du
faisceau $\mathscr{O}_X$

\oldpage[I]{169}
tel que $\Gamma(U,\mathscr{O}_X)=\bigcap_{x\in U}\mathscr{O}_x$ pour tout
ouvert $U\subset X$, avec une définition évidente des homomorphismes de
restriction. Nous laissons au lecteur le soin de vérifier qu'on obtient bien
ainsi un schéma intègre, dont les anneaux locaux sont les éléments de $X'$~;
nous n'utiliserons pas ce résultat par la suite.
\end{env}
